\subsection{The complete security statement and what remains conditional}\label{sec:mv:boundary}

\begin{definition}[Per-proof conditional false-parity bound]\label{def:mv:row-error}
Let $Q$ be an upper bound on the number
of affine parity rows tested during one
protocol execution. Assume every row
commitment and its extracted individual
bit openings are fixed before the
\emph{first inner bit challenge} of that row.
For any row whose parity is false
on these extracted bits, suppose the
probability that its $\kappa$ sequential
OR mini-proofs all accept, conditional
on the entire earlier visible history,
is at most $\eta_{\rm row}$.
In the stand-alone resettable
model with a uniform discrete-log
reduction bound $\delta_{\rm DL}$,
Theorem~\ref{thm:mv:false-parity-bound}
allows the choice
\[
 \eta_{\rm row}:=
 \min\{1,\,2^{-\kappa}+\sqrt{\delta_{\rm DL}}\}.
\]
Applying this bound uniformly over
adaptive histories is a hypothesis of
the composed security theorem.
\end{definition}

\begin{lemma}[Union bound for false accepted parity rows]\label{lem:mv:row-union}
Under Definition~\ref{def:mv:row-error},
the probability that any false
parity row is accepted is at most
$\min\{1,Q\eta_{\rm row}\}$.
\end{lemma}
\begin{proof}
Order the at most $Q$ row checks
by their occurrence in the protocol.
For each row, condition on its
prior visible history and the
already extracted bit openings.
If the row is false, its
conditional acceptance probability
is at most $\eta_{\rm row}$ by
hypothesis; if it is true, it
contributes nothing to the
false-row event. Averaging over
histories bounds each false-row
acceptance event by $\eta_{\rm row}$.
The probability union bound
over at most $Q$ rows gives
$Q\eta_{\rm row}$, capped by one.
\end{proof}

\begin{definition}[Global extraction error]\label{def:mv:coupling-error}
Write $\varepsilon_{\rm couple}$ for
a proven failure bound of an actual
multi-stage extractor satisfying
Definition~\ref{def:mv:global-coupling}.
Write $\varepsilon_{\rm oracle}$ for
a separately proven probability of
an accepted but inconsistent
external polynomial-commitment
or authenticated-opening relation.
These quantities are \emph{parameters}
and are not assigned numerical values
until the corresponding constructions
and reductions are supplied.
\end{definition}

\begin{theorem}[Conditional soundness of the sequential-bit compiler]\label{thm:mv:conditional-sound}
Assume the honest uniform outer
challenges, field representation,
packing consistency and fixed-value
conditions of
Theorem~\ref{thm:zkrs:false-sound}.
Suppose the implemented protocol
satisfies the global extraction
and oracle assumptions of
Definition~\ref{def:mv:coupling-error},
and each false accepted compiled
parity row obeys
Definition~\ref{def:mv:row-error}.
For every false public evaluation
claim,
\[
 \boxed{
 \Pr[\mathrm{accept}]
 \leq\min\left\{1,\,
 \frac{\ell+n}{2^n}
 +\varepsilon_{\rm couple}
 +Q\eta_{\rm row}
 +\varepsilon_{\rm oracle}
 \right\}.}
\]
\end{theorem}
\begin{proof}
Let $B_1$ be failure of the
global extraction coupling,
$B_2$ the event that a false
compiled parity row is accepted,
and $B_3$ an external oracle or
packing-link failure. Their
probabilities are bounded by
$\varepsilon_{\rm couple}$,
$Q\eta_{\rm row}$, and
$\varepsilon_{\rm oracle}$
by hypothesis and
Lemma~\ref{lem:mv:row-union}.
Outside $B_1\cup B_2\cup B_3$,
all registered field variables
are fixed at their required
outer challenge boundaries
and every accepted affine
constraint holds on the
extracted field values.
The coupled accepting run is
then an accepting run of the
ideal algebraic checker of
Chapter~9, with the same
independent uniform
$\theta,z$.
Theorem~\ref{thm:zkrs:false-sound}
bounds the probability of
accepting the false public claim
by $(\ell+n)/2^n$.
Add the probabilities of the
three excluded events and
cap by one.
\end{proof}

\begin{lemma}[Outer challenge assumptions]\label{lem:mv:outer-adaptive}
The malicious-verifier simulation
Theorem~\ref{thm:mv:affine-mvzk}
allows arbitrary causal outer
challenges $\theta,z$.
The soundness bound of
Theorem~\ref{thm:mv:conditional-sound}
requires these outer challenges
to be fresh independent uniforms.
\end{lemma}
\begin{proof}
The simulation proof proceeds by
comparing conditional transition
laws at each public history.
Lemmas~\ref{lem:zkrs:U-conditional}
and~\ref{lem:zkrs:A-conditional}
hold after fixing any causal
outer challenge values, and
the inner simulator matches
arbitrary verifier challenge bits.
By contrast, the algebraic
root bounds invoked in
Theorem~\ref{thm:zkrs:false-sound}
are Schwartz--Zippel bounds
that count roots under uniform
independent $\theta,z$.
Without that sampling condition
their numerical probabilities
do not follow.
\end{proof}

\begin{theorem}[An unmasked authenticated opening defeats witness privacy]\label{thm:mv:raw-oracle-leak}
Suppose an otherwise perfectly
simulated single-opening verifier
view additionally publishes the
entire unmasked partial vector
$s^*$, or publishes the full
transposed pad $c^p$ together
with the masked vector $m$.
Then there exist two admissible
witnesses with the same public
evaluation instance whose
resulting complete views are
perfectly distinguishable.
\end{theorem}
\begin{proof}
For the first disclosure rule,
use the two equal-evaluation
witnesses constructed in
Theorem~\ref{thm:dpr:not-zk}.
Their partial vectors differ
in one coordinate, which
directly distinguishes the
views with probability one.
For the second rule,
Theorem~\ref{thm:pack:reveal-leak}
gives the deterministic
reconstruction
$s^*=m+\mathcal T(c^p)$.
Apply the same distinguishing
test to the reconstructed
partial vector. All other
transcript messages can be
ignored by the distinguisher.
\end{proof}

\begin{proposition}[The Chapter 11 transcript result does not certify original KBFold]\label{prop:mv:not-kbfold}
Theorem~\ref{thm:mv:affine-mvzk}
does not imply an end-to-end
zero-knowledge or quantum-soundness
theorem for the original
odd-characteristic KBFold PCS.
\end{proposition}
\begin{proof}
The transcript modeled in
Definition~\ref{def:mv:affine-protocol}
contains only binary-field
masked Ring Switching messages,
Pedersen bit commitments and
their affine proof transcripts.
It does not contain KBFold's
odd-characteristic folded codeword
oracles, authenticated openings,
or its proximity and final-closure
checks. Moreover the simulator
uses classical rewinding of the
malicious verifier, which is not
the quantum simulation model.
The conclusion therefore has
a strictly smaller transcript
domain and different field and
adversary hypotheses than the
claimed PCS statement.
\end{proof}

\begin{remark}[What Chapter 11 genuinely advances]
The honest-sigma-verifier restriction
of Chapter~10 is removed for the
\emph{concrete affine transcript}
by sequential one-bit challenges
and a perfect expected-time
classical rewinding simulator.
Local temporal extraction is proved
with an exact two-run forking
probability. The global extraction
contract, malicious verifier access
to the underlying oracle and the
post-quantum KBFold composition
are still independent obligations.
\end{remark}
